\section{2-SAT}
\whenuse{Elegir verdadero/falso para cada variable bajo restricciones de dos
literales: ``al menos uno'', incompatibilidades e implicaciones.}

Cada variable tiene dos nodos: falso $2v$ y verdadero $2v+1$. La negación se
obtiene con XOR 1. Una cláusula $(a\lor b)$ equivale a las implicaciones
$(\neg a\Rightarrow b)$ y $(\neg b\Rightarrow a)$.

\complexity{$O(V+E)$ tiempo y memoria usando SCC sobre $2V$ literales.}

\begin{lstlisting}[style=cpp]
struct TwoSAT {
    int variables;
    vector<vector<int>> implication;

    explicit TwoSAT(int variable_count)
        : variables(variable_count),
          implication(2 * variable_count) {}

    // Nodo que representa (variable == value).
    int literal(int variable, bool value) const {
        return 2 * variable + (value ? 1 : 0);
    }

    void add_implication(int from, int to) {
        implication[from].push_back(to);
    }

    // (x == value_x) OR (y == value_y).
    void add_or(int x, bool value_x, int y, bool value_y) {
        int a = literal(x, value_x);
        int b = literal(y, value_y);
        add_implication(a ^ 1, b);
        add_implication(b ^ 1, a);
    }

    // Obliga a que variable == value.
    void force(int variable, bool value) {
        add_or(variable, value, variable, value);
    }

    // (x == value_x) implica (y == value_y).
    void imply(int x, bool value_x, int y, bool value_y) {
        add_or(x, !value_x, y, value_y);
    }

    // Exactamente una de las dos condiciones debe cumplirse.
    void add_xor(int x, bool value_x, int y, bool value_y) {
        add_or(x, value_x, y, value_y);
        add_or(x, !value_x, y, !value_y);
    }

    // x e y deben tener el mismo valor.
    void add_equal(int x, int y) {
        imply(x, true, y, true);
        imply(y, true, x, true);
    }

    // x e y deben tener valores diferentes.
    void add_different(int x, int y) {
        add_xor(x, true, y, true);
    }

    optional<vector<bool>> solve() const {
        vector<int> component =
            strongly_connected_components(implication);

        vector<bool> assignment(variables);
        for (int variable = 0; variable < variables; ++variable) {
            int is_false = literal(variable, false);
            int is_true = literal(variable, true);

            if (component[is_false] == component[is_true])
                return nullopt;

            assignment[variable] =
                component[is_true] > component[is_false];
        }
        return assignment;
    }
};
\end{lstlisting}

Operaciones frecuentes:

\begin{lstlisting}[style=cpp]
TwoSAT sat(3);
int A = 0, B = 1, C = 2;
sat.add_or(A, true, B, true);  // A o B
sat.force(C, true);            // C debe ser verdadero
sat.imply(A, true, C, true);   // A => C
sat.add_xor(A, true, B, true); // Exactamente uno: A o B
sat.add_equal(A, C);           // A == C
sat.add_different(B, C);       // B != C
auto answer = sat.solve();
if (!answer) { cout << "Imposible\n";
} else {
    for (bool value : *answer)
        cout << value << ' ';
    cout << '\n';
}
\end{lstlisting}

\warning{\texttt{strongly\_connected\_components} debe devolver las componentes
en orden topológico compatible con la comparación usada en \texttt{solve}.
Según la implementación de SCC, puede ser necesario cambiar
\texttt{>} por \texttt{<}. También se requiere \texttt{\#include <optional>}.}